% Lemma 3.1 of Smith, reordered to make the positional claim the first step.
\begin{frame}[t]{X1--X3: how the proof works}
  \begin{itemize}
    \item<1-> \textbf{(X1)} Every edge $\vec{f}$ contains exactly one or two switches.
    \item<2-> \textbf{(X2)} If $\vec f$ contains $x_i$ and $x_{i+1}$,
      then $x_{i+1}=f_4^-$.
    \item<3-> \textbf{(X3)} Every monochromatic edge in $c_X$ is a $(3,4)$-edge.
  \end{itemize}
  \medskip
  \begin{columns}[T,onlytextwidth]
    \begin{column}{0.48\textwidth}
      \centering
      \uncover<2->{
        \textbf{X2}\par\smallskip
        \begin{shiftchain}[
          column sep=6mm, row sep=6mm, dot radius=3.0pt,
          labels=false, label font=\scriptsize,
          bars={2,7}, bar labels={2/$x_i$,7/$x_{i+1}$},
        ]{8}
          \foreach \pos/\lab in {1/f_1,3/f_2,5/f_3,8/f_4}
            \node at (\pos,0) {$\lab$};
          \shiftedge{1/cfred,3/cfnavy,5/cfnavy,8/cfred}[$\vec f$]
        \end{shiftchain}
        \par\smallskip{\small $x_{i+1}=f_4^-$}
      }
    \end{column}
    \begin{column}{0.48\textwidth}
      \centering
      \uncover<3->{
        \textbf{X3}\par\smallskip
        \begin{shiftchain}[
          column sep=6mm, row sep=6mm, dot radius=3.0pt,
          labels=false, label font=\scriptsize,
          bars={5,7}, bar labels={5/$x_i$,7/$x_{i+1}$},
        ]{8}
          \foreach \pos/\lab in {1/f_1,3/f_2,5/f_3,8/f_4}
            \node at (\pos,0) {$\lab$};
          \shiftedge{1/cfred,3/cfred,5/cfred,8/cfred}[$\vec f$]
        \end{shiftchain}
        \par\smallskip{\small Both switches in $(f_3,f_4)$.}
      }
    \end{column}
  \end{columns}
\end{frame}
\begin{frame}{X1: every edge contains at least one switch}
  \medskip

  \begin{center}
    \only<1>{
    \begin{shiftchain}[
      column sep=6mm, row sep=5mm, dot radius=3.0pt,
      label font=\scriptsize,
    ]{12}

    \shiftvertexcolor{2,5,8,12}{cfred}
      \shiftedge{2,5,8,12}[$\vec f$]

    \end{shiftchain}
    }

    \only<2>{
    \begin{shiftchain}[
      column sep=6mm, row sep=5mm, dot radius=3.0pt,
      label font=\scriptsize,
      bars={11},
      bar labels{11/{$x$}}
    ]{12}

    \shiftvertexcolor{2,5,8}{cfred}
      \shiftedge{2,5,8,12/cfnavy}[$\vec f$]

    \end{shiftchain}
  }
  \end{center}
  At least 1 switch by greedy colouring design -- the greedy algorithm would always put a \(x=f_4^-\) switch
\end{frame}

\begin{frame}{X2: If there are at least two switches, the second is immediately before $f_4$}
  \begin{center}
  \begin{overlayarea}{\textwidth}{2.5cm}
    \centering
    \only<1>{
      \begin{shiftchain}[
        column sep=6mm, row sep=6mm, dot radius=3.0pt,
        label font=\scriptsize,
        bars={3}, bar labels={3/$x_i$},
      ]{12}
        \shiftedge{1/cfred,5/cfnavy,8/cfnavy,12/cfgrey}[$\vec f$]
      \end{shiftchain}
    }
    \only<2-3>{
      \begin{shiftchain}[
        column sep=6mm, row sep=6mm, dot radius=3.0pt,
        label font=\scriptsize,
        bars={3,10}, bar labels={3/$x_i$,10/$x_{i+1}$},
      ]{12}
        \shiftedge{1/cfred,5/cfnavy,8/cfnavy,12/cfred}[$\vec f$]
        \shiftedge{6/cfnavy,7/cfnavy,9/cfnavy,11/cfred}[$\vec e$]
      \end{shiftchain}
    }
    \only<4->{
      \begin{shiftchain}[
        column sep=6mm, row sep=6mm, dot radius=3.0pt,
        label font=\scriptsize,
        bars={3,10/dashed,11}, bar labels={3/$x_i$,11/$x_{i+1}$},
        ghost dots={11/2}, ghost dot style={cfred},
      ]{12}
        \shiftedge{1/cfred,5/cfnavy,8/cfnavy,12/cfred}[$\vec f$]
        \shiftedge{6/cfnavy,7/cfnavy,9/cfnavy,12/cfred}[$\vec e$]
      \end{shiftchain}
    }
  \end{overlayarea}
  \end{center}
  \begin{itemize}
    \item<1-> The first switch $x_i$ lies inside $\vec f$.
    \item<2-> There must be an edge $\vec e$ such that without $x_{i+1}$ $\vec e$ would be monochromatic. 
    \item<3-> $\vec e>x_i$, $x_{i+1}=e_4^-$
    \item<4-> If $e_1>f_1$ then if $e_4<f_4$, $\vec e$ and $\vec f$ are incomporable
    \item<5-> Thus $e_4=f_4=x_{i+1}$

  \end{itemize}

  \onslide<6->{\claimreminder{\Xone\Xtwo}}
\end{frame}

\begin{frame}[t]{}
  \small\setstretch{1.05}
  \claimreminder{\Xone\Xtwo\switchparity}

  \begin{lemma}{Claim}
    \Xthree
  \end{lemma}

  \begin{proof}{Proof}
    \begin{itemize}
      \item<1-> Let $\vec f$ be monochromatic in $c_X$.
      \item<2-> By X1 and Parity, it has exactly two switches in one gap.
      \item<3-> By X2, $x_{i+1}=f_4^-$ lies in $(f_3,f_4)$.
      \item<4-> Both switches lie in this gap, so $\vec f$ is a $(3,4)$-edge.
    \end{itemize}
  \end{proof}

  \begin{center}
    \begin{shiftchain}[
      column sep=6mm, row sep=6mm, dot radius=3.0pt,
      label font=\scriptsize,
      bars={7,11}, bar labels={7/$x_i$,11/$x_{i+1}$},
    ]{12}
      \shiftedge{2/cfred,4/cfred,6/cfred,12/cfred}[$\vec f$]
    \end{shiftchain}
  \end{center}
\end{frame}

\begin{frame}{Stage X: conclusions}
  \begin{itemize}
    \item \textbf{(X1)} Every edge $\vec f$ has exactly one or two switches.
    \item \textbf{(X2)} If $\vec f$ contains $x_i$ and $x_{i+1}$,
      then $x_{i+1}=f_4^-$.
    \item \textbf{(X3)} Every monochromatic edge in $c_X$ is a $(3,4)$-edge.
  \end{itemize}
\end{frame}
